% ═══════════════════════════════════════════════════════════════
% MUSTERLÖSUNG: Pretérito Indefinido (Einführung -ar Verben)
% ═══════════════════════════════════════════════════════════════
% Sequenz: 08 - Las vacaciones
% Stunde: 08c - Pretérito Indefinido Einführung
% ═══════════════════════════════════════════════════════════════

\documentclass[11pt, a4paper]{article}

% ═══════════════════════════════════════════════════════════════
% SW-MODUS TOGGLE
% ═══════════════════════════════════════════════════════════════
\newif\ifswmodus
\swmodusfalse

% ═══════════════════════════════════════════════════════════════
% PAKETE
% ═══════════════════════════════════════════════════════════════

\usepackage[utf8]{inputenc}
\usepackage[T1]{fontenc}
\usepackage[ngerman]{babel}
\usepackage{helvet}
\renewcommand{\familydefault}{\sfdefault}

\usepackage[margin=2cm]{geometry}
\usepackage{enumitem}
\usepackage{tabularx}
\usepackage{booktabs}
\usepackage{multirow}

\usepackage{xcolor}
\usepackage{framed}
\usepackage{tcolorbox}
\tcbuselibrary{skins}

\usepackage{fancyhdr}
\usepackage{lastpage}
\usepackage{needspace}

% ═══════════════════════════════════════════════════════════════
% FARBEN
% ═══════════════════════════════════════════════════════════════

\definecolor{spanisch-color}{HTML}{c41e3a}
\definecolor{grau-color}{HTML}{888888}
\definecolor{hellgrau-color}{HTML}{f5f5f5}
\definecolor{orange-color}{HTML}{b35c00}
\definecolor{loesung-color}{HTML}{c62828}

\definecolor{sw-dark}{HTML}{000000}
\definecolor{sw-gray}{HTML}{666666}
\definecolor{sw-white}{HTML}{ffffff}

\ifswmodus
    \colorlet{spanisch}{sw-dark}
    \colorlet{grau}{sw-gray}
    \colorlet{hellgrau}{sw-white}
    \colorlet{orange}{sw-dark}
    \colorlet{loesung}{sw-dark}
\else
    \colorlet{spanisch}{spanisch-color}
    \colorlet{grau}{grau-color}
    \colorlet{hellgrau}{hellgrau-color}
    \colorlet{orange}{orange-color}
    \colorlet{loesung}{loesung-color}
\fi

\newcommand{\fachfarbe}{spanisch}

% ═══════════════════════════════════════════════════════════════
% VARIABLEN
% ═══════════════════════════════════════════════════════════════

\newcommand{\thema}{Pretérito Indefinido (Einführung -ar)}
\newcommand{\sequenz}{08}
\newcommand{\block}{c}
\newcommand{\fach}{Spanisch}
\newcommand{\klasse}{7}
\newcommand{\maxpunkte}{24}
\newcommand{\datum}{\today}

% ═══════════════════════════════════════════════════════════════
% KOPF- UND FUSSZEILE
% ═══════════════════════════════════════════════════════════════

\pagestyle{fancy}
\fancyhf{}
\renewcommand{\headrulewidth}{0.4pt}
\renewcommand{\footrulewidth}{0.4pt}

\lhead{\fach{} -- Klasse \klasse}
\chead{\textcolor{loesung}{\textbf{ML-\sequenz\block{} (MUSTERLÖSUNG)}}}
\rhead{\datum}

\lfoot{}
\cfoot{}
\rfoot{Seite \thepage\ von \pageref{LastPage}}

% ═══════════════════════════════════════════════════════════════
% BEFEHLE
% ═══════════════════════════════════════════════════════════════

\newcommand{\punktebox}[1]{\hfill\fbox{\small \_/#1}}

\newcommand{\loesung}[1]{\textcolor{loesung}{\textbf{#1}}}

\newcommand{\luecke}[1]{\rule{#1}{0.4pt}}

\newtcolorbox{merkkasten}[1][]{
    colback=hellgrau,
    colframe=\fachfarbe,
    colbacktitle=white,
    coltitle=\fachfarbe,
    boxrule=1pt,
    arc=0pt,
    left=8pt, right=8pt, top=8pt, bottom=8pt,
    fonttitle=\bfseries,
    attach boxed title to top left={yshift=-2mm, xshift=5mm},
    boxed title style={boxrule=0pt, colback=white},
    title=#1
}

\newtcolorbox{vokabelkasten}{
    colback=white,
    colframe=\fachfarbe,
    boxrule=0.5pt,
    arc=0pt,
    left=5pt, right=5pt, top=5pt, bottom=5pt
}

\newtcolorbox{hinweis}{
    colback=white,
    colframe=orange,
    colbacktitle=white,
    coltitle=orange,
    boxrule=1pt,
    arc=0pt,
    left=5pt, right=5pt, top=5pt, bottom=5pt,
    fonttitle=\bfseries,
    attach boxed title to top left={yshift=-2mm, xshift=5mm},
    boxed title style={boxrule=0pt, colback=white},
    title=Tipp
}

\newenvironment{aufgabe}[1]{%
    \needspace{5\baselineskip}%
    \nopagebreak[4]%
    \vspace{10pt}
    \noindent\textbf{\textcolor{\fachfarbe}{Aufgabe #1}}
    \nopagebreak[4]%
    \vspace{5pt}
    \nopagebreak[4]%
    \begin{list}{}{\leftmargin=0pt}
    \item[]
}{%
    \end{list}
}

\newcommand{\es}[1]{\textcolor{\fachfarbe}{\textbf{#1}}}

% ═══════════════════════════════════════════════════════════════
% DOKUMENT
% ═══════════════════════════════════════════════════════════════

\begin{document}

% ═══════════════════════════════════════════════════════════════
% KOPFBEREICH
% ═══════════════════════════════════════════════════════════════

\begin{center}
    {\Large\bfseries\textcolor{\fachfarbe}{Arbeitsblatt: \thema}}\\[0.3em]
    {\large\textcolor{loesung}{\textbf{--- MUSTERLÖSUNG ---}}}
\end{center}

\vspace{5pt}

\noindent
\begin{tabularx}{\textwidth}{|X|X|X|}
\hline
\textbf{Name:} \textcolor{loesung}{Musterschüler/in} &
\textbf{Klasse:} \klasse &
\textbf{Datum:} \datum \\
\hline
\end{tabularx}

\vspace{15pt}

% ═══════════════════════════════════════════════════════════════
% MERKKASTEN: Konjugation (mit Lösungen)
% ═══════════════════════════════════════════════════════════════

\begin{merkkasten}[Pretérito Indefinido: Konjugation -ar Verben]
\textbf{Das Pretérito Indefinido} beschreibt \textbf{abgeschlossene} Handlungen in der Vergangenheit.

\vspace{0.5em}

\begin{center}
\begin{tabular}{|l|c|l|}
\hline
\textbf{Person} & \textbf{Endung} & \textbf{hablar (sprechen)} \\
\hline
yo & -é & habl\loesung{é} \\
tú & -aste & habl\loesung{aste} \\
él/ella/usted & -ó & habl\loesung{ó} \\
nosotros/as & -amos & habl\loesung{amos} \\
vosotros/as & -asteis & habl\loesung{asteis} \\
ellos/ellas/ustedes & -aron & habl\loesung{aron} \\
\hline
\end{tabular}
\end{center}

\end{merkkasten}

\vspace{10pt}

% ═══════════════════════════════════════════════════════════════
% MERKKASTEN: Signalwörter (mit Lösungen)
% ═══════════════════════════════════════════════════════════════

\begin{merkkasten}[Signalwörter (palabras clave)]
Diese Wörter zeigen: Es geht um die Vergangenheit!

\vspace{0.5em}

\begin{tabular}{ll}
\es{ayer} & = \loesung{gestern} \\[0.8em]
\es{el año pasado} & = \loesung{letztes Jahr} \\[0.8em]
\es{el verano pasado} & = \loesung{letzten Sommer} \\[0.8em]
\es{la semana pasada} & = \loesung{letzte Woche} \\
\end{tabular}
\end{merkkasten}

\vspace{15pt}

% ═══════════════════════════════════════════════════════════════
% AUFGABE 1: Signalwörter zuordnen (mit Lösungen)
% ═══════════════════════════════════════════════════════════════

\begin{aufgabe}{1}
Verbinde die spanischen Signalwörter mit der deutschen Bedeutung. \punktebox{4}
\end{aufgabe}

\begin{center}
\begin{tabular}{rcl}
\es{ayer} & \loesung{$\rightarrow$} & \loesung{gestern} \\[0.8em]
\es{el año pasado} & \loesung{$\rightarrow$} & \loesung{letztes Jahr} \\[0.8em]
\es{la semana pasada} & \loesung{$\rightarrow$} & \loesung{letzte Woche} \\[0.8em]
\es{el verano pasado} & \loesung{$\rightarrow$} & \loesung{letzten Sommer} \\
\end{tabular}
\end{center}

\textit{\textcolor{grau}{Bewertung: 1 Punkt pro richtiger Zuordnung}}

\vspace{15pt}

% ═══════════════════════════════════════════════════════════════
% AUFGABE 2: Verbformen bilden (mit Lösungen)
% ═══════════════════════════════════════════════════════════════

\begin{aufgabe}{2}
Konjugiere die Verben im Pretérito Indefinido. \punktebox{6}
\end{aufgabe}

\begin{vokabelkasten}
\textbf{Verben:} hablar (sprechen) | viajar (reisen) | nadar (schwimmen) | visitar (besuchen)
\end{vokabelkasten}

\vspace{0.5em}

\noindent
a) yo / hablar: \loesung{hablé} \\[0.8em]
b) tú / viajar: \loesung{viajaste} \\[0.8em]
c) él / nadar: \loesung{nadó} \\[0.8em]
d) nosotros / visitar: \loesung{visitamos} \\[0.8em]
e) vosotros / hablar: \loesung{hablasteis} \\[0.8em]
f) ellos / viajar: \loesung{viajaron} \\

\textit{\textcolor{grau}{Bewertung: 1 Punkt pro richtiger Form. Akzent erforderlich bei -é und -ó!}}

\vspace{15pt}

% ═══════════════════════════════════════════════════════════════
% AUFGABE 3: Lückentext (mit Lösungen)
% ═══════════════════════════════════════════════════════════════

\begin{aufgabe}{3}
Ergänze die Lücken mit der richtigen Form des Pretérito Indefinido. \punktebox{6}
\end{aufgabe}

\begin{hinweis}
Achte auf die Akzente: -é und -ó haben einen Akzent!
\end{hinweis}

\vspace{0.5em}

\noindent
\textbf{Mis vacaciones en Mallorca}

\vspace{0.5em}

\noindent
El verano pasado mi familia y yo \loesung{viajamos} (viajar) a Mallorca.

\vspace{0.5em}

\noindent
Yo \loesung{nadé} (nadar) mucho en el mar.

\vspace{0.5em}

\noindent
Mi hermana \loesung{tomó} (tomar) el sol en la playa.

\vspace{0.5em}

\noindent
Nosotros \loesung{visitamos} (visitar) la catedral de Palma.

\vspace{0.5em}

\noindent
Mis padres \loesung{hablaron} (hablar) español con los camareros.

\vspace{0.5em}

\noindent
¿Y tú? ¿\loesung{Viajaste} (viajar) el verano pasado?

\textit{\textcolor{grau}{Bewertung: 1 Punkt pro richtiger Form}}

\vspace{15pt}

% ═══════════════════════════════════════════════════════════════
% AUFGABE 4: Eigener Urlaubsbericht (Beispiellösung)
% ═══════════════════════════════════════════════════════════════

\begin{aufgabe}{4}
Schreibe mindestens 4 Sätze über deinen letzten Urlaub (oder erfinde einen Urlaub). \punktebox{8}
\end{aufgabe}

\textbf{Beispiellösung:}

\loesung{El año pasado viajé a Italia con mi familia.}

\loesung{Yo nadé en el mar Mediterráneo.}

\loesung{Mi madre tomó el sol en la playa.}

\loesung{Nosotros visitamos Roma y el Coliseo.}

\loesung{Mi padre habló italiano con los camareros.}

\loesung{¡Las vacaciones fueron muy divertidas!}

\vspace{0.5em}

\textit{\textcolor{grau}{Bewertungskriterien:}}
\begin{itemize}
    \item Mindestens 4 Sätze vorhanden: 2P
    \item Signalwort verwendet: 1P
    \item Verschiedene Personen: 2P
    \item Verschiedene Verben: 2P
    \item Sprachliche Korrektheit (Konjugation + Akzente): 1P
\end{itemize}

% ═══════════════════════════════════════════════════════════════
% PUNKTEÜBERSICHT
% ═══════════════════════════════════════════════════════════════

\vspace{20pt}
\hrule
\vspace{10pt}

\begin{flushright}
\textbf{Punkte:} \loesung{\maxpunkte} / \maxpunkte
\end{flushright}

\end{document}
